Put $\theta=\angle \dfrac{A}{2}$ and $s=\sin\theta$. The circle tangent to $AB$ and $AC$ at points three units from $A$ has centre $J$ on the angle bisector, with $AJ=\dfrac{3}{\cos\theta}$ and radius $3\tan\theta$. In the isosceles triangle, the incenter satisfies $AI=\dfrac{7\cos\theta}{(1+s)}$. Since $I$ lies on this circle and on its axis, $AI=AJ\pm3\tan\theta$. The minus sign would give $7(1-s^2)=3(1-s^2)$, impossible. The plus sign gives $7(1-s)=3(1+s)$, so $s=\dfrac{2}{5}$. Consequently $BC=14s=\boxed{\dfrac{28}{5}}$.
\begin{asy}
import graph;import geometry;usepackage("amsmath");size(7cm);
real h=7*sqrt(21)/5;pair A=(0,h),B=(-2.8,0),C=(2.8,0),J=A+(0,-15/sqrt(21)),I=(0,2*sqrt(21)/5),D=A+3/7*(B-A),E=A+3/7*(C-A);draw(A--B--C--cycle);draw(circle(J,6/sqrt(21)),blue);draw(A--J,dashed);dot(I);label("$A$",A,(0,1));label("$B$",B,(-1,-1));label("$C$",C,(1,-1));label("$D$",D,(-1,0));label("$E$",E,(1,0));label("$I$",I,(1,0));label("$J$",J,(-1,0));
\end{asy}

